% Generated draft: requires mathematical and source review.
\documentclass[11pt]{article}
\usepackage[margin=25mm]{geometry}
\usepackage{fontspec}
\setmainfont{Latin Modern Roman}
\usepackage{amsmath,amssymb,mathtools,unicode-math}
\setmathfont{Latin Modern Math}
\usepackage{enumitem,needspace,xurl,hyperref}
\setlength{\parindent}{0pt}
\setlength{\parskip}{5pt}
\setlength{\emergencystretch}{3em}
\newcommand{\N}{\mathbb{N}}
\newcommand{\Z}{\mathbb{Z}}
\newcommand{\Q}{\mathbb{Q}}
\newcommand{\R}{\mathbb{R}}
\newcommand{\C}{\mathbb{C}}
\newcommand{\F}{\mathbb{F}}
\newcommand{\A}{\mathbb{A}}
\newcommand{\PP}{\mathbb{P}}
\newcommand{\E}{\mathbb{E}}
\newcommand{\eps}{\varepsilon}
\newcommand{\dd}{\,\mathrm{d}}
\begin{document}
\section*{UnsolvedMath ID: 30006445}
\textbf{Source problem code:} catalogue entry
\subsection{Definitions and mathematical statement}
A log scheme is a scheme \(Y\) equipped with a sheaf of commutative monoids \(\mathcal M_Y\) and a homomorphism \(\alpha:\mathcal M_Y\to\mathcal O_Y\), where \(\mathcal O_Y\) is viewed multiplicatively, such that \(\alpha^{-1}(\mathcal O_Y^\times)\to\mathcal O_Y^\times\) is an isomorphism.

Let \(\mathcal C_{\mathrm{src}}\) denote the class of log schemes obtained by the construction in part (b) of the supplied proposition, within that proposition's intended setup. Its inputs include a polytope \(P\), its dual \(P^\vee\), and a compatible collection \((f_\Delta)_{\Delta}\) of zero-mutable Laurent polynomials, one for each facet \(\Delta\) of \(P^\vee\). The supplied proposition asserts that such a collection gives a global generically log smooth structure \(\mathcal M_Y\) on \(Y\). Part (a) asserts that an individual choice on \(\Delta\) gives a singular generically log smooth structure at the one-dimensional stratum of \(Y\) corresponding to \(\Delta\).

\textbf{Conjecture (attributed in the source to Corti-R.).} For every \((Y,\mathcal M_Y)\in\mathcal C_{\mathrm{src}}\), there exists a log-crepant log resolution of \((Y,\mathcal M_Y)\), and there also exists a smoothing of \((Y,\mathcal M_Y)\), in the senses intended by the source. These are two existence assertions; the supplied wording does not impose a compatibility condition between the resolution and the smoothing.

\textbf{Unresolved scope and terminology.} The supplied extract does not specify the full construction of \(Y\), the definition of zero-mutability, the compatibility condition on the facet polynomials, or the precise meanings of log-crepant log resolution and smoothing in this setting. It also omits the ground field, log base, ambient dimension for the conjecture, and hypotheses on the log structures. Thus \(\mathcal C_{\mathrm{src}}\) is a source-dependent class, and the conjecture above is not yet an intrinsic, fully specified formulation. In particular, no smoothing base, condition on its fibers, or formula for log crepancy is imposed here without source evidence.

\textbf{Supplied contextual theorem.} The extract separately states that every one of the 105 smooth Fano threefolds can be obtained by smoothing a toroidal crossing Fano variety \(Y\) arising from a central subdivision of a reflexive polytope \(P\), using for its log structure only products of the facet polynomials denoted \(f_{\mathrm{Tom}}\), \(f_{\mathrm{Jerry}}\), and \(f_{A_n}\) on facets of \(P^\vee\). These polynomial names are retained from the source; their definitions and the range of \(n\) are not supplied. This contextual assertion does not supply the missing general hypotheses of the conjecture or assert its universal resolution conclusion.
\subsection{Short English statement}
The conjecture attributed to Corti-R. asks whether every log scheme produced by the source's compatible facetwise zero-mutable Laurent polynomial construction admits both a log-crepant log resolution and a smoothing. The construction and several essential terms are incompletely specified in the supplied extract, so the statement remains conditional on the original setup. The supplied partially solved label and the contextual assertion about 105 smooth Fano threefolds are source reports requiring independent mathematical review.
\subsection{Sources}
\begin{itemize}
\item Local UnsolvedMath registry: supplied complete source record, permanent ID \texttt{30006445}, title \emph{Crepant Resolutions and Smoothings of Toroidal Fano Varieties}, exact source problem code \texttt{catalogue entry}. The conjecture is attributed in the extract to ``Corti-R.''; no expansion of that attribution is supplied. The record gives proposed year 2025.
\item Supplied primary citation: \emph{Singularities (2026)}, Oberwolfach Reports, DOI 10.4171/owr/2025/45, pp.\ 2413--2490. The literature-source label is ``Oberwolfach Report 45/2025.'' \url{https://doi.org/10.4171/owr/2025/45}. The extract does not give a conjecture-specific page or a complete account of the preceding setup.
\item Supplied literature source: \emph{How to make log structures}. \url{https://www.cambridge.org/core/journals/moduli/article/how-to-make-log-structures/974D2003CB24AD0B7DCD4D7507A011E6}.
\item Supplied literature source: \emph{Smoothings from zero-mutable Laurent polynomials}. \url{https://arxiv.org/abs/2503.18661}.
\item Supplied literature assessment, checked at the exact recorded date \texttt{2026-08-22}: status ``partially solved.'' The curation says that the universal conjecture is not established, that the report proves the smoothing statement for all 105 smooth Fano threefolds in its constructed class, and that later work constructs log-crepant resolutions and smoothings in additional families. This is supplied dated curation derived from the OWR record, not an independently verified literature assessment. No source was visited for this draft.
\end{itemize}
\end{document}
